% Source-audited historical subsection for a future JoA manuscript revision.
% Keep it distinct from the V0--V2 evaluator results and the later V15 screens.
\subsection{The Wright aircraft as an evolving aerodynamic and control system}
\label{sec:wright-history-dynamics}

The 1903 Flyer is an observed flight baseline, not a certificate of passive
stability or a fixed template for the three model-generated concepts. The
Wrights' route began with kites and successive piloted gliders; shortfalls in
the 1901 lift predictions led them to measure wing shapes in their own wind
tunnel, revise the lifting surfaces, and refine three-axis control in the 1902
glider. The powered Flyer retained two fabric-covered, braced wing decks,
wing-warp roll control, a forward pitch-control surface, and a linked aft
vertical rudder. A prone pilot lay beside the engine on the lower deck, while
chains drove two counter-rotating pusher propellers. Thus a propeller or
center-of-gravity change altered an integrated aircraft rather than an
isolated subsystem \cite{ref2,ref35,ref6}.

The original structure matters aerodynamically. Wood spars, ribs, fabric,
struts, and tensioned wires supplied the load path while permitting the
intended wing warp. In a later full-scale replica test, fabric billowing
increased effective camber and changed the lift curve, and the flexible
canard showed substantially different effectiveness from smaller rigid
models. That replica was locally reinforced to accept a tunnel sting, so its
measurements are valuable but are not proof of the exact 1903 airframe's
strength or an as-flown dynamic derivative database \cite{ref13}. A proposed
historical-style aircraft must therefore report member and joint stiffness,
control-surface travel under load, and fabric/shape uncertainty, not just a
three-view outline.

Propulsion was likewise part of the balance problem. The 1903 machine used a
roughly 12-hp engine, two pusher propellers and a chain transmission; the
propellers were treated as rotating lifting surfaces, not simple flat
blades \cite{ref35,ref43}. In the 1999 full-scale replica tunnel program,
the reproduced propellers at the test stand's permitted maximum of 340 rpm
did not produce thrust equal to drag at the selected 25-knot condition.
That limitation of a particular powered replica test does not overturn the
recorded 1903 flights. It demonstrates why engine rpm, delivered shaft power,
propeller advance ratio, transmission losses and installed drag must be
matched before declaring any of the present concepts capable of level
powered flight \cite{ref13}. Historical notebook static-thrust values and
modern replica thrust measurements also refer to different apparatus and
conditions; they cannot be transplanted to a model-designed propeller
\cite{ref43}.

The subsequent Wright machines were revisions, not repetitions of Flyer I.
The 1904 Flyer II kept the broad canard-biplane arrangement but changed
camber, structure and power plant. Moving mass aft initially aggravated
pitch handling; forward ballast and later control-arm changes followed.
The 1905 Flyer III then underwent several changes in one season, including
removal of forward vertical ``blinkers,'' an enlarged canard, enlarged and
more distant aft vertical surfaces, and a transition toward more independent
rudder operation. Its longer control moment arms increased authority and
pitch damping even though the canard configuration did not become passively
stable in pitch. Treating an early 1905 short-nose configuration and the
autumn final machine as the same aircraft would erase the causal design
cycle \cite{ref41,ref42,ref10}.

In the 1905 study, the engine carried over from 1904 is reported at roughly
21 hp by that stage, driving two counter-rotating propellers; the source
does not provide an installed performance map for every 1905 revision
\cite{ref10}. A change in power or propeller design must therefore be
evaluated alongside altered wing loading, drag, center of gravity and
control moments, rather than treated as an independent ``more power'' fix.

Later canard Wright machines also cannot
be assigned the 1903 coefficients: the 1909 configuration is reported as
more pitch-unstable than 1903 in a modern reconstruction, and the move to
an aft horizontal tail around 1910 finally changed the passive-stability
architecture \cite{ref41}. No common, source-verified derivative table for
every intervening variant is claimed here.

The later production path illustrates why a recognizable Wright silhouette
is not a sufficient comparator. The 1909 Military Flyer still used a front
elevator, wing warping, an aft rudder, and two chain-driven pusher
propellers, but carried two people and had a smaller wing area and higher
speed than earlier machines \cite{ref44}. The 1910 Model B moved pitch
control to an aft elevator and dispensed with the front canard
\cite{ref41,ref45}. Thus a ``Wright-like'' claim must identify the year,
loading, geometry, control linkage, and powertrain; the 1903 and 1910
aircraft do not even share the same longitudinal-control topology.

The cited studies answer different questions and should not be merged into
one purported ``Wright dataset.'' Culick and Jex analyzed the 1903
aircraft's aerodynamics, instability and pilot-in-the-loop control; their
work is a flight-mechanics reconstruction, not a direct flight-test record
of every derivative \cite{ref6}. Jex and colleagues obtained full-scale,
six-component wind-tunnel evidence from a carefully constructed but
reinforced 1903 replica, revealing fabric and canard-flexibility effects
and an unfavorable propulsion result at the tested rpm \cite{ref13}.
Padfield and Lawrence reconstructed the 1904--1905 modification cycle and
pilot handling, while Lawrence and Padfield combined 1905 tunnel data,
computed rate derivatives and piloted simulation to compare 1903 and final
1905 modes \cite{ref42,ref10}. Papachristodoulou and Culick extended the
configuration-level flight-mechanics comparison into the later canard and
aft-tail years \cite{ref41}. Ash and colleagues independently examined
historical propeller shapes and static-thrust tests; static thrust is not
an installed thrust--power--speed map \cite{ref43}. These papers justify
the kinds of evidence sought here, but none supplies the absent engine,
structure or unsteady-flow data for Astra, Fable or Opus.

\subsection{Measured derivatives, reconstructed modes, and pilot action}
\label{sec:wright-mode-baseline}

To compare signs without conflating conventions, define the conventional
fixed-control static margin as
$SM=-C_{m_\alpha}/C_{L_\alpha}$, with the derivatives evaluated at a stated
trim point and with the same reference chord. The 1905 study instead reports
$H_n=\partial C_M/\partial C_L$ as a positive \emph{magnitude of
instability}; its $H_n$ must not be read as a positive restoring margin.
Define the nondimensional pitch-rate derivative by
$C_{m_q}=\partial C_m/\partial(q\bar c/2V)$ and the corresponding
unsteady incidence-rate derivative by
$C_{m_{\dot\alpha}}=\partial C_m/\partial(\dot\alpha\bar c/2V)$.
The historical sources inspected here support selected measured slopes and
modelled rate terms, but do not provide a verified common
$C_{m_{\dot\alpha}}$ for all Wright configurations \cite{ref13,ref10,ref42}.

\begin{table}[p]
\centering\small
\caption{Selected published Wright stability quantities. Values belong to the stated replica or simulation and must not be substituted into a V15 design.}
\label{tab:wright-derivative-baseline}
\begin{tabularx}{\linewidth}{@{}p{.19\linewidth}p{.21\linewidth}p{.22\linewidth}X@{}}
\toprule
Configuration & Quantity & Reported value & Evidence and qualification\\
\midrule
1903 Flyer replica & Conventional $SM$ & about $-0.25$ ($-25\%$ MAC) & Full-scale Ames tunnel interpretation; reference CG near $0.30\bar c$, neutral point near $0.05\bar c$; approximately 25 kt \cite{ref13}.\\
1903 Flyer model & $C_{m_q}$ & about $-1.53$ & Estimated pitch damping in the Lawrence--Padfield model, not a direct 1903 in-flight measurement \cite{ref10}.\\
1903 Flyer model & unstable pitch pole; doubling time & $+1.6861\ \mathrm{s}^{-1}$; $0.41$ s & Linearized reconstruction; varies with mass and flight condition \cite{ref10}.\\
1905 final Flyer III & $C_{n_\beta}$ & $+0.0403\ \mathrm{rad}^{-1}$ & Reported final-configuration directional slope over $-16^\circ<\beta<16^\circ$ in the model/tunnel study \cite{ref10}.\\
1905 final Flyer III & $C_{m_q}$ & about $-2.62$; alternative $-4.67$ & Two published computational estimates; disagreement is retained, not averaged \cite{ref10,ref41}.\\
1905 final Flyer III & unstable pitch pole; doubling time & $+0.3162\ \mathrm{s}^{-1}$; $2.192$ s & Linearized reconstruction with slower divergence than the 1903 model despite greater static instability \cite{ref10}.\\
1905 final Flyer III & roll; Dutch roll; spiral poles & $-3.8984$; $-0.4944\pm1.2038i$; $+0.2499\ \mathrm{s}^{-1}$ & Linearized 26-kt model; the positive spiral pole remains divergent \cite{ref10}.\\
\bottomrule
\end{tabularx}
\end{table}

The 1905 analysis also gives a reported 1903 directional slope
$C_{n_\beta}\simeq+0.0368\ \mathrm{rad}^{-1}$ from the earlier Ames work;
this is a secondary transcription, not a value independently re-extracted
from the original tunnel plots \cite{ref10,ref13}. An earlier 2005 account
prints $+0.00403\ \mathrm{rad}^{-1}$ for the final 1905
$C_{n_\beta}$, whereas the later detailed study prints
$+0.0403\ \mathrm{rad}^{-1}$. Until the original figure's slope and
normalization are reconciled, numerical cross-paper comparison must flag
this factor-of-ten discrepancy rather than silently choose a value
\cite{ref42,ref10}. The intermediate 1905 directional derivatives plotted
for blinkers, canard and vertical-tail revisions are authors' estimates;
only the specified final test configuration should be called directly
measured \cite{ref10}.

These results explain how the Wrights could fly without modern eigenvalue
analysis. A positive open-loop pitch pole means a small disturbance grows
unless an external action intervenes; it does not mean that no controllable
trajectory exists. The pilot observed pitch attitude and moved the canard
continuously. In the reconstructed pilot--aircraft model, increasing
feedback could suppress the nonoscillatory pitch divergence but could also
reduce damping of another oscillatory mode, especially when perception and
muscle delay were represented. The same pitch-attitude feedback that was
tolerable in the more damped 1905 model destabilized an oscillatory mode of
the 1903 model; this is a model-based explanation of high workload, not a
reconstructed command history of a particular historical pilot
\cite{ref6,ref10,ref42}. For lateral control, wing warp produced adverse
yaw; the linked rudder helped initial correction, but tight low-speed turns
could combine sideslip, inboard-wing high incidence and a rapidly growing
spiral. The 1905 analysis attributes improved recovery to changed
rudder authority/linkage and greater control of speed during turns
\cite{ref10,ref42}. These observations are historical analysis, not
instructions to operate an untested replica.

The implication for the present experiment is methodological. For Fable V15,
the independent, power-off wing--tail AVL sweep yields a local balance at
$V=13$ m/s under assumed mass and section data and a window-sensitive
\emph{apparent} margin near $0.53$--$1.23\%$ of the reference chord.
This is not comparable in fidelity to a six-component, propeller-on Ames
replica test, and it does not supply $C_{m_{\dot\alpha}}$, validated
installed $C_{m_q}$, lateral--directional derivatives or modal roots.
Astra and Opus V15 retain independent propulsion or installation holds;
none of the three designs currently has a verified whole-aircraft trim or
accepted dynamic mode set. Our target is therefore first a physically
specified and independently checked airframe/propulsion model, next a
six-force-and-moment trim over a declared envelope, then traceable
derivatives, poles and sensitivity, and only then a separate analysis of
pilot or augmentation authority. A Wright-like aircraft need not possess
passive stability in every open-loop mode to be historically flyable, but
without both the missing model evidence and a control-authority test we
cannot claim that any V15 aircraft would be flyable by a pilot.

The proposed designs also diverge from the Wright machines in ways that
matter to the calculation. Astra's tandem surfaces distribute lifting and
pitch-control duties differently from a dominant Wright biplane plus
canard; a chosen center of gravity or lift split cannot substitute for a
measured six-component balance. Fable's biplane and twin pusher layout
looks closer to the 1903 architecture, yet its assumed engine--propeller
installation, firewall and tank separation, loaded structural clearances,
and high-drag power case remain open. Opus proposes another pusher
arrangement, but its predicted thrust does not close its stated drag case.
For all three, the structural load path must preserve intended control
travel, and a feasible pilot-control mapping requires signed control
derivatives, available travel and rate, and delay as well as free-aircraft
modal roots. These are explicit acceptance tests, not properties conferred
by a familiar shape or by copying a historical pilot response.
